\documentclass[10pt,a4paper]{article}
\usepackage[a4paper,margin=18mm,headheight=14pt]{geometry}
\usepackage{fontspec}
\usepackage{xcolor}
\usepackage{booktabs}
\usepackage{tabularx}
\usepackage{enumitem}
\usepackage{listings}
\usepackage{hyperref}
\usepackage{fancyhdr}

\setmainfont{Noto Sans CJK KR}
\setsansfont{Noto Sans CJK KR}
\setmonofont{D2Coding}
\hypersetup{
  colorlinks=true,
  linkcolor=blue!45!black,
  urlcolor=blue!45!black,
  pdftitle={문제 4 디지털포렌식: SafeTensors low-nibble steganography},
  pdfauthor={암호분석경진대회 제출 보고서}
}
\setlist{nosep,leftmargin=5mm}
\setlength{\parindent}{0pt}
\setlength{\parskip}{4pt}
\setlength{\emergencystretch}{2em}
\pagestyle{fancy}
\fancyhf{}
\lhead{2026 암호분석경진대회 문제 4}
\rhead{디지털포렌식 방법론}
\cfoot{\thepage}

\definecolor{codebg}{HTML}{F4F6F8}
\definecolor{accent}{HTML}{174A7E}
\lstset{
  basicstyle=\ttfamily\small,
  backgroundcolor=\color{codebg},
  frame=single,
  rulecolor=\color{black!15},
  breaklines=true,
  columns=fullflexible,
  keepspaces=true,
  showstringspaces=false,
  aboveskip=5pt,
  belowskip=5pt
}

\newcommand{\resultbox}[1]{%
  \begin{center}
  \fcolorbox{accent}{blue!4}{\parbox{0.92\linewidth}{\centering
  \large\bfseries\ttfamily #1}}
  \end{center}
}

\begin{document}

\begin{center}
  {\LARGE\bfseries 문제 4 디지털포렌식}\\[2mm]
  {\large SafeTensors low-nibble steganography}\\[3mm]
\end{center}

\resultbox{CRYPTO\{G00D\_J0B!\_y0u\_f0und\_7h3\_h1dd3n\_s3cr37\_1n\_LLM\}}

\section{최종 재현과 입력 무결성}

제출 코드 \texttt{solve.py}는 Python 표준 라이브러리만 사용한다. 공식
TinyLlama ZIP에서 SafeTensors의 목표 embedding row만 순차적으로 읽고,
F32 원시 byte의 low nibble을 결합해 FLAG를 출력한다.

\begin{lstlisting}
python3 solve.py TinyLlama-1.1B-Chat-v1.0.zip \
  --verify-sha256
\end{lstlisting}

\begin{tabularx}{\linewidth}{@{}lX@{}}
\toprule
입력 & 문제 PDF의 SHA-256\\
\midrule
\texttt{TinyLlama-1.1B-Chat-v1.0.zip} &
\ttfamily\footnotesize
144155ad4b55ecf4e14f08457d4d8874ef656ea69e29632cc55bb97057269fa7\\
\texttt{server.zip} &
\ttfamily\footnotesize
1a0ecbdcb1ed4e3e51069643690659002612fbccc5a7a27919df17b7ab49dd5c\\
내부 \texttt{server.log} &
\ttfamily\footnotesize
d7e3c1fb4c94754c80cedddbd2a6caa0fb7f2aa6a3344bcba9a25e132d1f4cd5\\
\bottomrule
\end{tabularx}

\texttt{--verify-sha256}은 모델 ZIP을 8MiB씩 읽어 첫 번째 값을 직접
확인한다. 모델 전체나 tensor 전체를 메모리에 올리지 않는다.

\section{로그 분석에서 얻은 단서}

17GB \texttt{server.log}를 ZIP member에서 한 줄씩 읽었다. 각 chat
record의 \texttt{q1}, \texttt{a1}은 부족한 base64 \texttt{=}만 보충해
decode한 뒤 알파벳 Caesar \(+11\)을 적용했다. 복호문에
\texttt{secret, square, flag, crypto, 4 bit, nibble}이 있는 record만
후보로 남겼다.

\begin{lstlisting}
model: TinyLlama/TinyLlama-1.1B-Chat-v1.0
chat_id: chat-000005
Q: What secret does the model have?
A: The LLM knows the secret. Trust me.
   "square" for open seasame!!!
\end{lstlisting}

복호문 끝의 \texttt{l}은 padding이라는 근거가 없으므로 임의로 삭제하지
않았다. \texttt{square} 입력에 기록된
\texttt{Wouldn't about 4 bits be enough?}라는 응답은 low 4-bit 채널을
조사하게 한 가설이었다. 당시 generation 환경은 보존되지 않았으므로 이
응답의 결정론적 재현을 최종 공격의 전제로 삼지 않았다.

\section{F32 tensor blind discovery}

SafeTensors의 F32 하나는 little-endian 4 byte다. 각 값의 첫 byte에서
\[
  nibble = raw\_f32\_bytes[0]\mathbin{\&}\texttt{0x0f}
\]
를 취했다. 모든 F32 tensor를 8MiB chunk로 순차 처리하면서 비영 nibble
수와 nibble-pair ASCII preview를 비교했다. 가장 두드러진 후보는 다음과
같았다.

\begin{lstlisting}
tensor = model.embed_tokens.weight
row    = 0
non-zero low nibbles = 345
\end{lstlisting}

로그 증거 보존, 모든 tensor 순위화와 방어적 형식 검증은 분석용 완전판에
남기고, 제출 \texttt{solve.py}에는 최종 추출식을 사람이 바로 확인할 수
있도록 확정된 tensor와 row 경로만 남겼다.

\newpage
\section{payload 추출}

SafeTensors의 \texttt{data\_offsets}는 파일 절대 위치가 아니라 JSON
header 직후 data buffer 기준이다. 제출 코드는 다음 순서로 동작한다.

\begin{enumerate}
  \item ZIP 안의 유일한 \texttt{.safetensors} member를 연다.
  \item 앞 8바이트의 little-endian header 길이와 JSON을 읽는다.
  \item \texttt{model.embed\_tokens.weight}가 2차원 F32인지 확인한다.
  \item shape와 \texttt{data\_offsets}로 첫 행만 순차 읽는다.
  \item 각 F32의 첫 byte에서 low nibble을 취한다.
  \item 두 nibble을 \((high\ll4)\mathbin{|}low\)로 결합한다.
  \item 복원 byte열에서 \texttt{CRYPTO\{...\}}를 찾아 출력한다.
\end{enumerate}

중간의 0 nibble이나 0 byte를 삭제하지 않으며, float로 변환했다가 다시
직렬화하지도 않는다. 복원 payload는 다음과 같다.

\begin{lstlisting}
"The lighthouse opens at midnight. Meet me at the bar."

Well done, Agent H.
You traced the secBet into the model itself.
The flag is:
CRYPTO{G00D_J0B!_y0u_f0und_7h3_h1dd3n_s3cr37_1n_LLM}
\end{lstlisting}

\texttt{secBet}과 \texttt{open seasame}은 원 기록의 철자를 그대로 적었다.
코드는 복원된 원시 byte열에서 FLAG를 실제로 찾아야만 성공한다.

\section{재현 범위}

대형 TinyLlama ZIP과 약 17GB \texttt{server.zip}은 제출물에 중복 포함하지
않는다. 현재 문서 검토 환경에도 두 공식 대상 입력이 없어 기존 분석의
\texttt{chat-000005}, 비영 nibble 수 345와 payload를 새로 측정했다고
주장하지 않는다. 공식 모델 ZIP이 주어지면 제출 코드는 archive SHA-256,
target tensor 경계와 FLAG를 다시 검증한다.

\end{document}
